%=====================================================================
%  FIRST-PAPER TEMPLATE — Mastering the Scientific Narrative
%  Resource Kit · C. G. Hounmenou (University of Labé / CERFIG)
%---------------------------------------------------------------------
%  HOW TO USE
%  1) Easiest: go to https://overleaf.com → New Project → Upload this
%     folder (main.tex + references.bib). Compiles out of the box.
%  2) For a REAL submission, download the OFFICIAL style of your venue
%     and paste your sections into it (never edit venue .sty files):
%       · NeurIPS  : https://neurips.cc  → "Paper Information / Style files"
%       · ICLR     : https://iclr.cc     → "Call for Papers / Author guide"
%       · ACL/AfricaNLP : https://github.com/acl-org/acl-style-files
%     Overleaf also hosts official templates: Menu → Templates → search
%     the venue name.
%  3) Every % comment below is a coach whispering. Delete them as you go.
%=====================================================================
\documentclass[11pt,a4paper]{article}

% ---------- Packages you will almost always need ----------
\usepackage[utf8]{inputenc}
\usepackage[T1]{fontenc}
\usepackage[margin=2.5cm]{geometry}
\usepackage{graphicx}      % figures
\usepackage{booktabs}      % professional tables (\toprule,\midrule,\bottomrule)
\usepackage{amsmath,amssymb,bm} % math; \bm{} for bold vectors
\usepackage{hyperref}      % clickable links & refs
\usepackage{xcolor}
\usepackage[numbers]{natbib} % citations: \citep{key} (Author, 2020) / \citet{key}

% ---------- Notation shortcuts: define ONCE, use everywhere ----------
% Consistency of notation is a top reviewer criterion (Checklist §2).
\newcommand{\x}{\bm{x}}          % vector: bold lowercase
\newcommand{\X}{\bm{X}}          % matrix: bold uppercase
\newcommand{\R}{\mathbb{R}}
\newcommand{\model}{AgroNet}     % <- rename once, updates everywhere

\title{%
% TITLE = your 15-word promise. Three formulas that work:
%   [Method]: [What it does] for [Problem]
%   [Question or claim as hook]?
%   [Resource]: a [scale] [type] for [context]
% Test: can a stranger guess method + problem + context from it alone?
\model{}: Lightweight Offline Detection of Cassava Diseases\\
for Low-End Mobile Devices}

\author{%
  First Author\thanks{Equal contribution / corresponding author: email@univ.edu} \\
  University of Labé, Guinea
  \and
  Second Author \\
  University of Lagos, Nigeria
% Your African affiliation is your lab entering the world map — state it
% proudly and completely (department, institution, country).
}
\date{}

\begin{document}
\maketitle

%=====================================================================
\begin{abstract}
% THE 5 MOVES (compress to 3 sentences if space is tight):
% [CONTEXT] stakes + numbers .......................................
Cassava feeds over 500 million Africans, yet mosaic and brown-streak
diseases destroy up to 40\% of yields.
% [GAP] what blocks today's solutions ..............................
Existing vision models for plant diagnosis require cloud connectivity
that is absent in most rural areas.
% [METHOD] one sentence, one system, one key property ..............
We introduce \model{}, an offline, quantized convolutional network
optimized for low-end mobile devices.
% [RESULTS] data + hard numbers, never adjectives ..................
Trained on 12{,}406 leaf images from 42 farms in Guinea and Nigeria,
\model{} reaches 78.3\% top-1 accuracy while cutting on-device latency
by 40\% versus MobileNetV3.
% [IMPACT] the door you open .......................................
We release our dataset, code and models to support agricultural AI
across the continent.
\end{abstract}

%=====================================================================
\section{Introduction}
% THE FUNNEL: general problem -> specific gap -> your contribution.
% Para 1 — CONTEXT: why the world should care (stakes, numbers, refs).
% Para 2 — GAP: what today's solutions cannot do, and why it matters
%          here (connectivity, compute, data coverage...). Cite fairly.
% Para 3 — CONTRIBUTION: "In this paper, we introduce ..." Answer the
%          Heilmeier questions: what? how is it done today? what is
%          new? who cares — locally AND globally?
% END WITH THE EXPLICIT LIST (reviewers hunt for it):
Our contributions are threefold:
\begin{itemize}
  \item \textbf{A new method:} we propose \model{}, the first \ldots
        designed for offline, low-compute settings (\S\ref{sec:method});
  \item \textbf{A new resource:} we release a dataset of 12{,}406
        images collected across 42 farms (\S\ref{sec:data});
  \item \textbf{New evidence:} \model{} outperforms strong baselines
        by 3.5 points while being $7\times$ faster (\S\ref{sec:results}).
\end{itemize}
% RULE: contributions are PROMISES; Results are PROOFS. Claim 3, prove 3.

%=====================================================================
\section{Related Work}
% Position yourself in the conversation — 3 mini-paragraphs, each:
%   "Prior line of work A did X \citep{...}. However, ..." -> your gap.
% Be generous and precise; the authors you cite may review you.

%=====================================================================
\section{Method}\label{sec:method}
% Gold standard: a competent reader could RE-IMPLEMENT from text alone.
% - Define every symbol at first use; keep the conventions above.
% - Figures > prose for architectures (make Figure 1 tell the story).
Let $\x \in \R^{d}$ denote an input image representation\ldots

%=====================================================================
\section{Dataset}\label{sec:data}
% The dataset biography (Checklist §2): provenance (where, when, by
% whom, consent/ethics), size & composition, preprocessing, splits
% (by FARM / by SPEAKER — no leakage), known biases, license.
% Publish a dataset card (Gebru et al., 2021) and archive with a DOI
% (Zenodo is free) — this is FAIR data \citep{wilkinson2016fair}.

%=====================================================================
\section{Experiments and Results}\label{sec:results}
% Protocol first: metrics (with formulas), splits, seeds, baselines.
% Report mean +/- std over >= 3 seeds. One figure = one message.

\begin{table}[t]
  \centering
  \caption{Top-1 accuracy (\%) on the cassava test set, mean$\pm$std
  over 10 seeds. Best in \textbf{bold}. Latency on a Tecno Spark 8.}
  \label{tab:main}
  \begin{tabular}{lcc}
    \toprule
    Model & Accuracy (\%) & Latency (ms) \\
    \midrule
    ResNet-50      & 71.2 $\pm$ 0.9 & 412 \\
    MobileNetV3    & 74.8 $\pm$ 0.8 & 96  \\
    EfficientNet-B0& 76.1 $\pm$ 0.7 & 198 \\
    \textbf{\model{} (ours)} & \textbf{78.3 $\pm$ 0.5} & \textbf{58} \\
    \bottomrule
  \end{tabular}
\end{table}

% Figures: export vector PDF from matplotlib (plt.savefig('fig.pdf'))
% \begin{figure}[t]\centering
%   \includegraphics[width=.8\linewidth]{figures/pareto.pdf}
%   \caption{Self-contained caption: what, where, and the takeaway.}
% \end{figure}

%=====================================================================
\section{Discussion and Limitations}
% Discussion: interpret — WHY it wins, WHEN it fails, what it means
% for deployment. Do not repeat the numbers.
% Limitations — the trust section; specific and honest, e.g.:
% "\model{} was validated on three cassava varieties; performance under
%  extreme low-light and on 2G-only devices remains untested."
% African-context constraints are valuable observations — own them.

%=====================================================================
\section{Conclusion}
% One paragraph: contribution + strongest number + the future work
% your limitations just wrote for you.

%=====================================================================
% Reproducibility: link an anonymized repo during review; public +
% Zenodo DOI at camera-ready. One command (run.sh) reproduces Table 1.

\bibliographystyle{plainnat}
\bibliography{references}

\end{document}
